\documentclass[12pt,a4paper]{article}
\usepackage[UTF8]{ctex}
\usepackage{amsmath, amssymb}
\usepackage{geometry}
\usepackage{booktabs}
\usepackage{float}
\geometry{left=2.5cm, right=2.5cm, top=2.5cm, bottom=2.5cm}
\title{\textbf{问题三 结果附表}\\[6pt]
\large 二分搜索判定过程 / d* 与 HPWL 汇总 / 临界点扰动核验}
\author{2026 年第七届"华数杯"大学生数学建模竞赛 B 题}
\date{2026/08/10}
\begin{document}
\maketitle

\section{三组芯片二分搜索可行性判定过程}

二分区间 $[0.075, 0.15]$（下界取 $d \ge 0.075$ 区间以聚焦临界点附近，
$0.15$ 为问题 2 基准已知可行），终止精度 $\varepsilon = 10^{-4}$；
\textbf{任一并行种子判定可行即可行}（OR 语义）。并行种子数
粗区间 5--7、精区间 10--15。

\subsection{n100（判定种子 15，全程 71.2s）}
\begin{table}[H]
\centering
\scriptsize
\setlength{\tabcolsep}{2.5pt}
\begin{tabular}{ccccccc}
\toprule
迭代 & 区间 $[lo, hi]$ & 测试点 $d$ & 轮廓 side & 判定 bbox & 可行 & 耗时(s) \\
\midrule
1 & $[0.075,0.15]$ & 0.075 & 440 & 270x1912 & 否 & 3.1 \\
2 & $[0.075,0.1125]$ & 0.1125 & 447 & 315x1945 & 是 & 2.6 \\
3 & $[0.075,0.0938]$ & 0.0938 & 444 & 295x1924 & 是 & 3.0 \\
4 & $[0.075,0.0844]$ & 0.0844 & 442 & 278x1888 & 是 & 3.2 \\
5 & $[0.075,0.0797]$ & 0.0797 & 441 & 224x1801 & 是 & 5.8 \\
6 & $[0.075,0.0773]$ & 0.0773 & 440 & 268x1990 & 是 & 6.0 \\
7 & $[0.075,0.0762]$ & 0.0762 & 440 & 306x1875 & 是 & 6.0 \\
8 & $[0.0756,0.0762]$ & 0.0756 & 440 & 356x1866 & 否 & 6.3 \\
9 & $[0.0759,0.0762]$ & 0.0759 & 440 & 378x1965 & 否 & 6.3 \\
10 & $[0.076,0.0762]$ & 0.076 & 440 & 309x1820 & 否 & 6.4 \\
11 & $[0.0761,0.0762]$ & 0.0761 & 440 & 279x1816 & 否 & 6.4 \\
\bottomrule
\end{tabular}
\end{table}

\subsection{n200（判定种子 15，全程 908.3s）}
\begin{table}[H]
\centering
\scriptsize
\setlength{\tabcolsep}{2.5pt}
\begin{tabular}{ccccccc}
\toprule
迭代 & 区间 $[lo, hi]$ & 测试点 $d$ & 轮廓 side & 判定 bbox & 可行 & 耗时(s) \\
\midrule
1 & $[0.075,0.15]$ & 0.075 & 435 & 249x2504 & 否 & 52.6 \\
2 & $[0.075,0.1125]$ & 0.1125 & 443 & 228x2591 & 是 & 50.4 \\
3 & $[0.075,0.0938]$ & 0.0938 & 439 & 254x2438 & 是 & 50.9 \\
4 & $[0.0844,0.0938]$ & 0.0844 & 437 & 264x2737 & 否 & 48.9 \\
5 & $[0.0844,0.0891]$ & 0.0891 & 438 & 262x2672 & 是 & 88.3 \\
6 & $[0.0844,0.0867]$ & 0.0867 & 437 & 247x2528 & 是 & 95.9 \\
7 & $[0.0844,0.0855]$ & 0.0855 & 437 & 260x2579 & 是 & 96.4 \\
8 & $[0.0844,0.085]$ & 0.085 & 437 & 261x2777 & 是 & 97.6 \\
9 & $[0.0847,0.085]$ & 0.0847 & 437 & 264x2511 & 否 & 97.4 \\
10 & $[0.0847,0.0848]$ & 0.0848 & 437 & 264x2466 & 是 & 94.6 \\
11 & $[0.0847,0.0847]$ & 0.0847 & 437 & 247x2597 & 是 & 90.9 \\
\bottomrule
\end{tabular}
\end{table}

\subsection{n300（判定种子 10，全程 2822.6s）}
\begin{table}[H]
\centering
\scriptsize
\setlength{\tabcolsep}{2.5pt}
\begin{tabular}{ccccccc}
\toprule
迭代 & 区间 $[lo, hi]$ & 测试点 $d$ & 轮廓 side & 判定 bbox & 可行 & 耗时(s) \\
\midrule
1 & $[0.075,0.15]$ & 0.075 & 542 & 266x3783 & 否 & 206.6 \\
2 & $[0.1125,0.15]$ & 0.1125 & 552 & 285x3863 & 否 & 208.6 \\
3 & $[0.1125,0.1312]$ & 0.1312 & 556 & 268x3961 & 是 & 203.6 \\
4 & $[0.1125,0.1219]$ & 0.1219 & 554 & 285x3771 & 是 & 208.1 \\
5 & $[0.1125,0.1172]$ & 0.1172 & 553 & 276x3560 & 是 & 257.3 \\
6 & $[0.1125,0.1148]$ & 0.1148 & 552 & 286x3904 & 是 & 290.8 \\
7 & $[0.1125,0.1137]$ & 0.1137 & 552 & 318x3850 & 是 & 256.8 \\
8 & $[0.1125,0.1131]$ & 0.1131 & 552 & 273x3920 & 是 & 280.6 \\
9 & $[0.1125,0.1128]$ & 0.1128 & 552 & 268x3796 & 是 & 292.8 \\
10 & $[0.1125,0.1126]$ & 0.1126 & 552 & 268x3791 & 是 & 256.7 \\
11 & $[0.1125,0.1126]$ & 0.1126 & 552 & 270x3834 & 是 & 245.7 \\
\bottomrule
\end{tabular}
\end{table}

\section{三组芯片最小死区比例与更新后 HPWL 汇总}

$d^*$ 处完整求解（多起点 10 轮取 HPWL 最优，仅接受可行轮次）得到
更新后 HPWL；全部经双向确认与独立合法性复核。

\begin{table}[H]
\centering
\scriptsize
\setlength{\tabcolsep}{2.5pt}
\begin{tabular}{lccccccc}
\toprule
芯片 & 判定种子 & $d^*$ & 轮廓 side & $d^*$ 处 HPWL（更新后） & 双向确认 & 冻结 $d^*$ & 较冻结 \\
\midrule

n100 & 15 & \textbf{0.0761} & 440 & 288184.0 & 通过 & 0.117307 & -35.2\% \\
n200 & 15 & \textbf{0.0841} & 437 & 559042.5 & 通过 & 0.108977 & -22.8\% \\
n300 & 10 & \textbf{0.1118} & 552 & 826375.0 & 通过 & 0.124885 & -10.5\% \\
\bottomrule
\end{tabular}
\end{table}

\section{最小死区比例附近扰动灵敏度核验}

双向确认在临界点附近做\textbf{上下扰动核验}：在 $d^*$ 处多种子判定
必须可行；在 $d^* - \varepsilon$（$\varepsilon = 10^{-4}$）处多种子
\textbf{必须全部不可行}——任一扰动方向上不满足则沿该方向移动
$\varepsilon$ 重验（最多 12 步），保证 $d^*$ 是"可行与不可行之间
验证过的最小稳定点"，而非二分收敛点的偶然结果。

\subsection{临界点扰动核验结果}

\begin{table}[H]
\centering
\scriptsize
\setlength{\tabcolsep}{2.5pt}
\begin{tabular}{lcccccc}
\toprule
芯片 & 判定种子 $K$ & $d^*$ & 核验点 $d^*$（可行种子数） & 核验点 $d^*-\varepsilon$（可行种子数） & 结论 \\
\midrule

n100 & 15 & \textbf{0.0761} & 可行（双向确认通过） & 不可行（0/15） & 验证通过 \\
n200 & 15 & \textbf{0.0841} & 可行（双向确认通过） & 不可行（0/15） & 验证通过 \\
n300 & 10 & \textbf{0.1118} & 可行（双向确认通过） & 不可行（0/10） & 验证通过 \\
\bottomrule
\end{tabular}
\end{table}

\subsection{确认阶段下降轨迹（二分上界 $\to$ $d^*$）}

每步：在 $d$ 处判定可行后，检查 $d - \varepsilon$；若仍可行则下移
$\varepsilon$ 重验，直至 $d - \varepsilon$ 处全部种子不可行。

\begin{table}[H]
\centering
\scriptsize
\setlength{\tabcolsep}{3pt}
\begin{tabular}{clcc}
\toprule
芯片 & 步 & 验证点 $d$ & 判定结果 \\
\midrule

n100 & 1 & 0.07617 & 可行；$d-\varepsilon=0.07607$ 处 1/15 种子仍可行（下移） \\
n100 & 2 & \textbf{0.07607} & 可行；$d-\varepsilon=0.07597$ 处 0/15 种子可行（下界确认通过） \\
n200 & 1 & 0.08474 & 可行；$d-\varepsilon=0.08464$ 处 1/15 种子仍可行（下移） \\
n200 & 2 & 0.08464 & 可行；$d-\varepsilon=0.08454$ 处 3/15 种子仍可行（下移） \\
n200 & 3 & 0.08454 & 可行；$d-\varepsilon=0.08444$ 处 1/15 种子仍可行（下移） \\
n200 & 4 & 0.08444 & 可行；$d-\varepsilon=0.08434$ 处 3/15 种子仍可行（下移） \\
n200 & 5 & 0.08434 & 可行；$d-\varepsilon=0.08424$ 处 3/15 种子仍可行（下移） \\
n200 & 6 & 0.08434 & 不可行（上移 $\varepsilon$ 重验） \\
n200 & 7 & 0.08434 & 可行；$d-\varepsilon=0.08424$ 处 2/15 种子仍可行（下移） \\
n200 & 8 & 0.08424 & 可行；$d-\varepsilon=0.08414$ 处 2/15 种子仍可行（下移） \\
n200 & 9 & 0.08414 & 可行；$d-\varepsilon=0.08404$ 处 1/15 种子仍可行（下移） \\
n200 & 10 & 0.08414 & 不可行（上移 $\varepsilon$ 重验） \\
n200 & 11 & \textbf{0.08414} & 可行；$d-\varepsilon=0.08404$ 处 0/15 种子可行（下界确认通过） \\
n300 & 1 & 0.11257 & 可行；$d-\varepsilon=0.11247$ 处 4/10 种子仍可行（下移） \\
n300 & 2 & 0.11247 & 可行；$d-\varepsilon=0.11237$ 处 1/10 种子仍可行（下移） \\
n300 & 3 & 0.11247 & 不可行（上移 $\varepsilon$ 重验） \\
n300 & 4 & 0.11247 & 可行；$d-\varepsilon=0.11237$ 处 4/10 种子仍可行（下移） \\
n300 & 5 & 0.11237 & 可行；$d-\varepsilon=0.11227$ 处 2/10 种子仍可行（下移） \\
n300 & 6 & 0.11227 & 可行；$d-\varepsilon=0.11217$ 处 2/10 种子仍可行（下移） \\
n300 & 7 & 0.11217 & 可行；$d-\varepsilon=0.11207$ 处 3/10 种子仍可行（下移） \\
n300 & 8 & 0.11207 & 可行；$d-\varepsilon=0.11197$ 处 4/10 种子仍可行（下移） \\
n300 & 9 & 0.11197 & 可行；$d-\varepsilon=0.11187$ 处 2/10 种子仍可行（下移） \\
n300 & 10 & 0.11187 & 可行；$d-\varepsilon=0.11177$ 处 2/10 种子仍可行（下移） \\
n300 & 11 & \textbf{0.11177} & 可行；$d-\varepsilon=0.11167$ 处 0/10 种子可行（下界确认通过） \\
\bottomrule
\end{tabular}
\end{table}

注：n100 从 $d=0.076172$ 下移 1 步至 $d^*=0.076072$；n200 从
$0.084741$ 经 10 步下降至 $0.084141$；n300 从 $0.112573$ 经 10 步
（含 1 次上移复核）下降至 $0.111773$。全部最终点满足
\texttt{d\_feas=True} 且 \texttt{below\_infeas=True}。

\end{document}
